\scanpage{62}
% archive sheet number: 61
on vérifie aisément que pour tt schém. abélien $B/S$, on
a une \uline{bijection}
\[
\mathrm{Hom}_{S\text{-gr}}(B,\underline{\mathrm{Pic}}_{X/S})\xleftarrow{\sim}\mathrm{Hom}_{S\text{-gr}}(B,\underline B_{X/S})
\]
d'où par \uline{dualité} \uline{des} V.A. une bijection
\[
\left\{\begin{aligned}
\mathrm{Hom}_{S\text{-gr}}(\widehat A,\underline{\mathrm{Pic}}_{X/S})&\xleftarrow{\sim}\mathrm{Hom}_{S\text{-gr}}(\underline A^0_{X/S},A)\\
\text{où}\qquad \underline A^0_{X/S}&=\widehat{\underline B}_{X/S}
\end{aligned}\right.
\]

Donc \uline{pour Un schém. abélien} $A$ \uline{sur} $S$, \deleted{\uline{\unsure{il existe|variable}}}
\uline{la donnée d'un fibré principal hom.} $P$ \uline{sous} $A$
\uline{et d'un $S$-morphisme} $\psi:X\longrightarrow P$ \uline{équivaut \unsure{can.|ess.} à}
\uline{la donnée d'un hom} \inserted{de \uline{$S$-gr}} $u:\underline A^0_{X/S}\longrightarrow A$. De façon
précise, il y a un fibré principal hom.
$\underline A^1_{X/S}$ sous $\underline A^0_{X/S}$, canonique, \deleted{tel que} \deleted{\unsure{tel que}}
\inserted{et un $S$-morphisme} \deleted{\unsure{unique}} $\psi_{X/S}:X\longrightarrow\underline A^1_{X/S}$, et la correspondance
\unsure{explicitée|explicite} plus haut s'obtient en prenant :
\[
\deleted{X} \quad P=\underline A^1_{X/S}\times_{\underline A^0_{X/S}}(A,u)\qquad \ldots\ldots\quad //
\]
On appelle
$\underline A^0_{X/S}$ et $\underline A^1_{X/S}$ resp. les schémas d'Albanese
\uline{homogènes} et \uline{inhomogènes} de $X/S$. Ils sont
en général non isom. \uline{On notera que}
$(\underline A^0_{X/S},\underline A^1_{X/S},\psi_{X/S})$ \uline{représente le foncteur} $T_X(P)$ par construction
